\noindent Abadie, A., A. Diamond, and J. Hainmueller. 2010. ``Synthetic Control Methods for Comparative Case Studies:  Estimating the Effect of California's Tobacco Control Program.'' \textit{Journal of the American Statistical Association} 105: 493-505.

\noindent Abadie, A., A. Diamond, and J. Hainmueller. 2015. ``Comparative Politics and the Synthetic Control Method.'' \textit{American Journal of Political Science} xx: yy-zz.

\noindent Abadie, A., and J. Gardeazabal. 2003.  ``The Economic Costs of Conflict:  A Case Study of the Basque Country.'' \textit{American Economic Review} 93: 113-132.

\noindent Acemoglu, D., S. Johnson, A. Kermani, J. Kwak, and T. Mitton.  2013.  ``The Value of Connections in the Turbulent Times:  Evidence from the United States.''  Working paper, at http://ssrn.com/abstract-2363609. 

\noindent Acemoglu, D., S. Johnson, and J., A. Robinson. ``The Colonial Origins of Comparative Development:  An Empirical Investigation.'',\textit{American Economic Review }91: 1369-1401.

\noindent Adams, R., and J. Santos. 2006. ``Identifying the Effect of Managerial Control on Firm Performance.''  \textit{Journal of Accounting and Economics }41: 55-85.

\noindent Adams, R., H. Almeida, and D. Ferreira. 2005. ``Powerful CEOs and Their Impact on Corporate Performance.''  \textit{Review of Financial Studies} 18: 1403-1432.

\noindent Akhigbe, A. and A. Martin. 2005. ``Valuation impact of Sarbanes--Oxley: Evidence from disclosure and governance within the financial services industry.'' \textit{Journal of Banking and Finance} 30: 989-1006.

\noindent Akhigbe, A., A. Martin, and M. Newman. 2010. ``Information Asymmetry Determinants of Sarbanes-Oxley Wealth Effects.'' \textit{Finance Management} 31: 1253-1272.

\noindent Aivazian, V., Y. Ge, and J. Qiu. 2005. ``Can Corporatization Improve the Performance of State-owned Enterprises even without Privatization?'' \textit{Journal of Corporate Finance} 11: 791-808.

\noindent Altamuro, J., and A. Beatty. 2010. ``How Does Internal Control Regulation Affect Financial Reporting?'' \textit{Journal of Accounting and Economics }49: 58-74.

\noindent Altamuro, J., A. Beatty, and J. Weber. 2005. ``The Effects of Accelerated Revenue Recognition on Earnings Management and Earnings Informativeness: Evidence from SEC Staff Accounting Bulletin No. 101.''  \textit{Accounting Review} 80: 373-401.

\noindent Altonji, J., T. Elder, and C. Taber. 2005. ``Selection on Observed and Unobserved Variables:  Assessing the Effectiveness of Catholic Schools.''  \textit{Journal of Political Economy} 113: 151-184.

\noindent Angrist, J., and W. Evans. 1998. ``Children and Their Parents' Labor Supply:  Evidence from Exogenous Variation in Family Size.'' \textit{American Economic Review} 88: 450-477.

\noindent Angrist, J., G. Imbens, and D. Rubin. 1996. ``Identification of Causal Effects Using Instrumental Variables.''  \textit{Journal of the American Statistical Association} 91: 444-455.

\noindent Angrist, J., and J-S. Pischke. 2009. ``Mostly Harmless Econometrics:  An Empiricist's Companion.''

\noindent Armstrong, C., A. Jagolinzer, and D. Larcker. 2010. ``Chief Executive Officer Equity Incentives and Accounting Irregularities.''  \textit{Journal of Accounting Research} 48: 225-271.

\noindent Atanasov, V., B. Black, C. Ciccotello, and S. Gyoshev. 2010. ``How Does Law Affect Finance?  An Examination of Equity Tunneling in Bulgaria.''  \textit{Journal of Financial Economics }96: 155-173.

\noindent Atanasov, V., and B. Black. 2016a. ``Database of Shocks Used in Corporate Finance and Accounting Research.'' working paper, at http://ssrn.com/abstract=2227717.

\noindent Atanasov, V., and B. Black. 2016b. ``The Trouble with Instruments:  Re-examining Shock-Based IV Designs.''\textit{ }working paper, at http://ssrn.com/abstract=2417689.

\noindent Atanasov, V., V. Ivanov, and K. Litvak. 2012. ``Does Reputation Limit Opportunistic Behavior in the VC Industry:  Evidence from Litigation against VDs.''  \textit{Journal of Finance} 67: 2215-2246. 

\noindent Autor, D. 2003. ``Outsourcing at Will:  The Contribution of Unjust Dismissal Doctrine to the Growth of Employment Outsourcing.''  \textit{Journal of Labor Economics} 21: 1-42.

\noindent Bae, K. and V. Goyal. 2010. ``Equity Market Liberalization and Corporate Governance.'' \textit{Journal of Corporate Finance} 609-621.

\noindent Baek, J., J. Kang, and K. Park. 2004. ``Corporate governance and firm value: evidence from the Korean financial crisis.''  \textit{Journal of Financial Economics} 265-313.

\noindent Bakke, T., and T. Whited. 2012. ``Threshold Events and Identification:  A Study of Cash Shortfalls.''  \textit{Journal of Finance} 67: 1083-1111

\noindent Bargeron, L., K. Lehn, and C. Zutter. 2010. ``Sarbanes-Oxley and Corporate Risk-Taking.''  \textit{Journal of Accounting and Economics} 49: 34-52.

\noindent Battalio, R., B. Hatch, and T. Loughran, 2011. `` Who benefited from the disclosure mandates of the 1964 Securities Act Amendments.'' \textit{Journal of Corporate Finance} 17: 1047-1063.

\noindent Bennedsen, M., K. Nielssen, F. Perez-Gonzalez, and D. Wolfenzon. 2007. ``Inside the Family Firm: The Role of Families in Succession Decisions and Performance.''  \textit{Quarterly Journal of Economics} 122: 647-691.

\noindent Berger, E., A. Butler, E. Hu, and M. Zekhnini.  2014.  ``Credit Be Dammed: The Impact of Banking Deregulation on Economic Growth.'' Working paper, at http://ssrn.com/abstract=2139679.

\noindent Berkman, H., R. Cole, and L. Fu. ``Political Connections and Minority-Shareholder, Protection: Evidence from Securities-Market Regulation in China.'' \textit{Journal of Financial and Quantitative Analysis} 1391-1417.

\noindent Bertrand, M., E. Duflo, and S. Mullainathan. 2004. ``How Much Should We Trust Differences-in-Differences Estimates?''  \textit{Quarterly Journal of Economics} 119: 249--75.

\noindent Bertrand, M., and S. Mullainathan. 2003. ``Enjoying the Quiet Life?  Corporate Governance and Managerial Preferences.''  \textit{Journal of Political Economy }111: 1043-1075.

\noindent Besley, T., and A. Case. 2000. ``Unnatural Experiments?  Estimating the Incidence of Endogenous Policies.''  \textit{Economic Journal} 110: F672-F694.

\noindent Bhagat, S., and R. Romano. 2002a. ``Event Studies and the Law: Part I -- Technique and Corporate Litigation.'' \textit{American Law and Economics Review} 4: 141-167.

\noindent Bhagat, S., and R. Romano. 2002b. ``Event Studies and the Law: Part II -- Empirical Studies of Corporate Law.''  \textit{American Law and Economics Review} 4: 380-423.

\noindent Black, B., A. de Carvalho, V. Khanna, W. Kim and B. Yurtoglu. 2014. ``Methods for Multicountry Studies of Corporate Governance and Evidence from the BRIKT Countries.'' \textit{Journal of Econometrics} 183: 230-240.

\noindent Black, B., H. Jang, and W. Kim. 2006. ``Does Corporate Governance Affect Firms' Market Values?  Evidence from Korea.''  \textit{Journal of Law, Economics and Organization} 22: 366-413.

\noindent Black, B., and W. Kim. 2012. ``The Effect of Board Structure on Firm Value:~ A Multiple Identification Strategy Approach Using Korean Data.''  \textit{Journal of Financial Economics }103: 203-226.

\noindent Black, B., W. Kim, H. Jang, and K. Park. 2013. ``How Corporate Governance Affects Firm Value:  Evidence from Korea on a Related Party Transactions Channel.'' working paper, at http://ssrn.com/abstract=844744.

\noindent Black, B. and V. Khanna. 2007. ``Can Corporate Governance Reforms Increase Firm Market Values? Event Study Evidence from India.''  \textit{Journal of Empirical Legal Studies} 4: 749-796.

\noindent Bowen, D., L. Fresard, and J. Taillard. 2015. ``What's Your Identification Strategy?  Innovation in Corporate Finance Research,'' working paper, at http://ssrn.com/abstract=2378460. 

\noindent Brochet, F. 2010. ``Information Content of Insider Trades before and after the Sarbanes-Oxley Act.''  \textit{The Accounting Review} 85: 419-446. 

\noindent Broughman, B., and J. Fried. 2010. ``Renegotiation of Cash Flow Rights in the Sale of VC-backed Firms.''  \textit{Journal of Financial Economics} 95: 384-399.

\noindent Brown, J., N. Liang, and S. Weisbenner. ``Executive Financial Incentives and Payout Policy: Firm Responses to the 2003 Dividend Tax Cut.''  \textit{Journal of Finance} 62: 1935-1965.

\noindent Brown, P., W. Beekes, and P. Verhoeven. 2011. ``Corporate Governance, Accounting and Finance:  A Review.''  \textit{Accounting and Finance} 51: 96-172.

\noindent Busso, M., J. DiNardo, and J. McCrary. 2014. ``New Evidence on the Finite Sample Properties of Propensity Score Reweighting and Matching Estimators.'' \textit{Review of Economics and Statistics }96:  885-897.

\noindent Byard, D., Y. Li, and Y. Yi. 2011. ``The Effect of Mandatory IFRS Adoption on Financial Analysts' Information Environment.''  No. 1 \textit{Journal of Accounting Research }49: 69-96.

\noindent Cameron, C., J. Gelbach, and D. Miller. 2008. ``Bootstrap-Based Improvements for Inference with Clustered Errors.''  \textit{Review of Economics and Statistics} 90: 414-427.

\noindent Cameron, C., J. Gelbach, and D. Miller. 2011. ``Robust Inference with Multiway Clustering.''  \textit{Journal of Business and Economic Statistics} 29: 238-249.

\noindent Catan, E., and M. Kahan. 2015.  ``The Law and Finance of Anti-Takeover Statutes.'' working paper, at http://ssrn.com/abstract=2517594.

\noindent Chang, H., J. Chen, W. Liao. 2006. ``Birendra K. Mishra, CEOs'/CFOs' Swearing by the Numbers: Does It Impact Share Price of the Firm?''  \textit{The Accounting Review} 81: 1-27.

\noindent Carvalhal da Silva, A., and A. Subrahmanyam. 2007. ``Dual-class Premium, Corporate Governance, and the Mandatory Bid Rule:  Evidence from the Brazilian Stock Market.''  \textit{Journal of Corporate Finance} 13: 1-24.

\noindent Chava, S., and A. Purnanandam. 2010. ``CEOs versus CFOs: Incentives and corporate policies.''\textit{ Journal of Financial Economics }97: 263-278.\textit{}

\noindent Chen, S., R. Chou, and S. Chou. 2009. ``The Impact of Investment Opportunities and Free Cash Flow on Financial Liberalization: A Cross-Firm Analysis of Emerging Economies.'' \textit{Financial Management} 543-566.

\noindent Cheng, Q., and D. Farber. 2008. ``Earnings Restatements, Changes in CEO Compensation, and Firm Performance.''  \textit{The Accounting Review} 83: 1217-1250

\noindent Chetty, R., and E. Saez. 2005. ``Dividend Taxes and Corporate Behavior: Evidence from the 2003 Dividend Tax Cut.''  \textit{Quarterly Journal of Economics} 120: 791-833.

\noindent Chhaochharia, V., and Y. Grinstein. 2007. ``Corporate Governance and Firm Value: The Impact of the 2002 Governance Rules.''  \textit{Journal of Finance} 62: 1789-1825. 

\noindent Chhaochharia, V., and Y. Grinstein. 2009. ``CEO Compensation and Board Structure.''  \textit{Journal of Finance} 64: 231-261.

\noindent Choi, J., S. Park, and S. Yoo,. 2007. ``The Value of Outside Directors:  Evidence from Corporate Governance Reform in Korea.''  \textit{Journal of Financial and Quantitative Analysis} 42: 941-962.

\noindent Cicero, D. 2009. ``The Manipulation of Executive Stock Option Exercise Strategies: Information Timing and Backdating.''  \textit{Journal of Finance }64:\textit{ }2627-2663.

\noindent Claessens, S., and B. Yurtoglu. 2012. ``Corporate Governance in Emerging Markets:  A Survey.'' working paper, at http://ssrn.com/abstract=1988880.

\noindent Coates, J. 2000.  ``Takeover Defenses in the Shadow of the Pill:  A Critique of the Scientific Evidence.''  \textit{Texas Law Review} 79: 271-382.

\noindent Cochran, W. 1965. ``The Planning of Observational Studies of Human Populations (with discussion).''  \textit{Journal of the Royal Statistical Society A} 128: 234--265.

\noindent Conley, T., C. Hansen, and P. Rossi. 2012. ``Plausibly Exogenous.'' \textit{Review of Economics and }Statistics 94: 260-272. 

\noindent Conley, T., and C. Taber. 2011. ``Inference with ``Difference in Differences'' with a Small Number of Policy Changes.''  \textit{Review of Economics and Statistics} 93: 113-125.

\noindent Cremers, M., and A. Ferrell. 2014. ``Thirty Years of Shareholder Rights and Firm Valuation.'' \textit{Journal of Finance }(forthcoming), working paper at http://ssrn.com/abstract=1413133.

\noindent Crump, R., V. Hotz, G. Imbens, and O. Mitnik. 2009. ``Dealing with Limited Overlap in Estimation of Average Treatment Effects''  \textit{Biometrika} 96: 187-199.

\noindent Dahya, J., J. McConnell, and N. Travlos. 2002. ``The Cadbury Committee, Corporate Performance, and Top Management Turnover.'' \textit{Journal of Finance} 57: 461-483.

\noindent Dahya, J., and J. McConnell. 2007. ``Board Composition, Corporate Performance, and the Cadbury Committee Recommendation.'' \textit{Journal of Financial and Quantitative Analysis }42: 535-564.

\noindent Daske, H., L. Hail, C. Leuz, and R. Verdi. 2008. ``Mandatory IFRS Reporting around the World: Early Evidence on the Economic Consequences.'' \textit{Journal of Accounting Research}, 46: 1085-1141.

\noindent De Jong,  A., D. DeJong, G. Mertensa, and C. Wasleyc. 2005. ``The Role of Self-Regulation in Corporate Governance: Evidence and Implications from the Netherlands.'' \textit{Journal of Finance} 11: 473-503. 

\noindent DeFond, M., X. Hu, M. Hung, and S. Li. 2011. ``The Impact of Mandatory IFRS Adoption on Foreign Mutual Fund Ownership: The Role of Comparability.''  \textit{Journal of Accounting and Economics} 51: 240-258.

\noindent Demirguc-Kunt, A., I. Love, and V. Maksimovic. 2006. ``Business Environment and the Incorporation Decision.''  \textit{Journal of Banking and Finance }30: 2967-2993.

\noindent Desai, M., and D. Dharmapala. 2009. ``Corporate Tax Avoidance and Firm Value.''  \textit{Review of Economics and Statistics }91: 537-546.

\noindent Dharmapala, D., F. Foley, and K. Forbes. 2011. ``Watch What I Do, Not What I Say.'' The Unintended Consequences of the Homeland Investment Act.''  \textit{Journal of Finance} 66: 753-787.

\noindent Dharmapala, D., and V. Khanna. 2013. ``Corporate Governance, Enforcement, and Firm Value:  Evidence from India.''  \textit{Journal of Law, Economics and Organization} 29: 1056-1082.

\noindent Dinc, S. 2005. ``Politicians and Banks:  Political Influences on Government-Owned Banks in Emerging Markets.''  \textit{Journal of Financial Economics} 77: 453-479.

\noindent Doidge, C., G. Karolyi, and R. Stulz. 2009. ``Has New York Become Less Competitive than London in Global Markets?  Evaluating Foreign Listing Choices Over Time.''  \textit{Journal of Financial Economics} 91: 253-277.

\noindent Doidge, C., G. Karolyi, and R. Stulz. 2010. ``Why Do Foreign Firms Leave U.S. Equity Markets?'' \textit{Journal of Finance} 65: 1507-1553.

\noindent Duchin, R., J. Matsusaka, and O. Ozbas. 2010. ``When Are Outside Directors Effective?'' \textit{Journal of Financial Economics} 95: 195-214.

\noindent Dunning, T. 2012. ``Natural Experiments in the Social Sciences:  A Design-Based Approach (Cambridge Univ. Press).''

\noindent Duflo, E., R. Hanna, and S. Ryan. 2010. ``Incentives Work:  Getting Teachers to Come to School.'' working paper.

\noindent Engel, E., R. Hayes, and X. Wang. 2010. ``Audit Committee Compensation and the Demand for Monitoring of the Financial Reporting Process.'' \textit{Journal of Accounting and Economics} 49: 136-154.

\noindent Engel, E., R. Hayes, and X. Wang. 2007. ``The Sarbanes--Oxley Act and Firms' Going-Private Decisions.'' \textit{Journal of Accounting and Economics} 44: 116-145.

\noindent Faccio, M. and D. Parsley. 2009. ``Sudden Deaths: Taking Stock of Geographic Ties.'' \textit{Journal of} \textit{Financial and Quantitative Analysis} 44: 683-718. 

\noindent Fahlenbrach, R., A. Low, and R. Stulz. 2010. ``Why Do Firms Appoint CEOs as Outside Directors.'' \textit{Journal of Financial Economics} 97: 12-32.

\noindent Fama, E., L. Fisher, M. Jensen, and R. Roll.  1969.  ``The Adjustment of Stock Prices to New Information.''  \textit{International Economic Review} 10:  1-21. 

\noindent Fernandes, N., U. Lel, and D. Miller. 2010. ``Escape from New York: The market impact of loosening disclosure requirements.''  \textit{Journal of Financial Economics} 95: 129-147.

\noindent Fich, E., and A. Shivdasani. 2007. ``Financial fraud, director reputation, and shareholder wealth.'' \textit{Journal of Financial Economics} 86: 306-336. 

\noindent Francis, B., I. Hasan, K. John, and L. Song. 2011. ``Corporate Governance and Dividend Payout Policy: A Test Using Antitakeover Legislation.''  \textit{Financial Management}, 37: 83-112

\noindent Francis, B., I. Hasan, K. John, and M. Waisman. 2010. ``The Effect of State Antitakeover Laws on the Firm's Bond Holders.'' \textit{Journal of Financial Economics }96: 127-154.

\noindent Frangakis, C., and D. Rubin. 2002. ``Principal Stratification in Causal Inference.'' \textit{Biometrics} 58: 21-29.

\noindent French, K., and R. Roll.  1986.  ``Stock Return Variances:  The Arrival of Information and the Reaction of Traders.''  \textit{Journal of Financial Economics} 17: 5-26.

\noindent Frumento, P., F. Mealli, B. Pacini, and D. Rubin. 2012. ``Evaluating the Effect of Training on Wages in the Presence of Noncompliance, Nonemployment, and Missing Outcome Data.'' \textit{Journal of the American Statistical Association} 107: 450-466.

\noindent Gagnon, L., and G. Karolyi. 2012. ``The Economic Consequences of the U.S. Supreme Court's \textit{Morrison v. National Australia Bank }Decision for Foreign Stocks Cross-Listed in U.S. Markets.'' working paper, at http://ssrn.com/abstract=1961178.

\noindent Galiania, S., A. Murphy, and J. Pantano. 2012. ``Estimating Neighborhood Choice Models:  Lessons from the Moving to Opportunity Experiment.'' working paper, at http://ssrn.com/abstract=2015474. 

\noindent Gao, F., J. Wu, and J. Zimmerman. 2009. ``Unintended Consequences of Granting Small Firms Exemptions from Securities Regulation: Evidence from the Sarbanes-Oxley Act.''  \textit{Journal of Accounting Research} 47: 459-506.

\noindent Gassen, J. 2014.  ``Causal Inference in Empirical Archival Financial Accounting Research.'' \textit{Accounting, Organizations and Society }39: 535-544.

\noindent Gelbach, J., E. Helland, and J. Klick. 2013. ``Valid Inference in Single-Firm, Single-Event Studies.''  \textit{American Law and Economics Review} 15: 495-541.

\noindent Gerber, A., and D. Green.  2012.  \textit{Field Experiments:  Design, Analysis and Interpretation}.

\noindent Ghosh, C., J. Harding, and B. Phani. 2008. ``Does Liberalization Reduce Agency Costs? Evidence from the Indian Banking Sector.'' \textit{Journal of Banking and Finance} 32: 405-419.

\noindent Giannetti, M., and L. Laeven. 2009. ``Pension Reform, Ownership Structure and Corporate Governance:  Evidence from a Natural Experiment.''  \textit{Review of Financial Studies} 22: 4092-4127.

\noindent Gilson, R., and B. Black. 1995. \textit{The Law and Finance of Corporate Acquisitions}, second edition.

\noindent Gippel, J., T. Smith, and Y Zhu.  2015.  ``Endogeneity in Accounting and Finance Research:  Natural Experiments as a State-of-the-Art Solution.''  \textit{Abacus} 51:  143-168. 

\noindent Giroud, X., and H. Mueller. 2010. ``Does Corporate Governance Matter in Competitive Industries?'' \textit{Journal of Financial Economics} (2010) 95: 312--331.

\noindent Glaeser, E., R. La Porta, F. Lopez-de-Silanes, and A. Shleifer. 2010. ``Do Institutions Cause Growth?'' \textit{Journal of Economic Growth} 9: 271-303.

\noindent Glaeser, E. 2008. ``Research Incentives and Empirical Methods, in Andrew Caplin and Andrew Schotter, eds.'' \textit{The Foundations of Positive and Normative Economics} (Oxford Univ. Press), at 300-319 (working paper at http://ssrn.com/abstract=934557).

\noindent Goldman, E., J. Rocholl, and J. So. 2009. ``Do Politically Connected Boards Affect Firm Value?'' \textit{Review of Financial Studies} 22: 2331-2360. 

\noindent Gomes, A., G. Gorton, and L. Madureira. 2007. ``SEC Regulation Fair Disclosure, Information, and the Cost of Capital.'' \textit{Journal of Corporate Finance} 113: 300-334.

\noindent Gompers, P., J. Ishii, and A. Metrick. 2003. ``Corporate Governance and Equity Prices.'' \textit{Quarterly Journal of Economics} 118: 107-155.

\noindent Greenstone, M., P. Oyer, and A. Vissing-Jorgensen. 2006. ``Mandated Disclosure, Stock Returns, and the 1964 Securities Acts Amendments.'' \textit{ Quarterly Journal of Economics} 121: 399-460.

\noindent Grullon, G., S. Michenaud, and J. Weston. 2012. ``The Real Effects of Stock Market Prices.'' working paper, at http://ssrn.com/abstract=1959615.

\noindent Guner, B., U. Malmendier, and J. Tate. 2008. ``Financial Expertise of Directors.'' \textit{Journal of Financial Economics }88: 323-354.

\noindent Hail, L., and C. Leuz. 2009. ``Cost of Capital Effects and Changes in Growth Expectations around U.S. Cross-Listings.''  \textit{Journal of Financial Economics} 93: 428-454.

\noindent Harvey, C., Y. Liu, and H. Zhu. 2013. . ``. . and the Cross-Section of Expected Returns.'' working paper, at http://ssrn.com/abstract=2249314.

\noindent Heckman, J., H. Ichimura, and P. Todd. 1997. ``Matching as an Econometric Evaluation Estimator:  Evidence from Evaluating a Job Training Programme.''  \textit{Review of Economic Studies} 64: 605-654.

\noindent Hennessy, C. A., and I. A. Strebulaev.  2015.  ``Beyond Random Assignment:  Credible Inference of Causal Effects in Dynamic Economies.''  Working paper, at http://ssrn.com/abstract=2564828.

\noindent Heron, Randall, and Eric Lie.  2015.  ``The Effect of Poison Pill Adoptions and Court Rulings on Firm Entrenchment.''  \textit{Journal of Corporate Finance}, forthcoming.

\noindent Hirano, K., and G. Imbens. 2004. ``The Propensity Score with Continuous Treatments, in Andrew Gelman and X.-L. Meng eds.'' \textit{Applied Bayesian Modeling and Causal Inference from Incomplete-Data Perspectives }73--84.

\noindent Ho, D., and K. Imai. 2005. ``Randomization Inference with Natural Experiments:  An Analysis of Ballot Effects in the 2003 California Recall Election.'' \textit{Journal of the American Statistical Association} 101: 888-900.

\noindent Hochberg, Y., and L. Lindsey. 2010. ``Incentives, Targeting and Firm Performance:  An Analysis of Non-executive Stock Options.''  \textit{Review of Financial Studies} 23: 4148-4186.

\noindent Hochberg, Y., P. Sapienza, and A. Vissing-Jorgensen. 2009. ``A Lobbying Approach to Evaluating the Sarbanes-Oxley Act of 2002.''  \textit{Journal of Accounting Research} 47: 519-583.

\noindent Holger, D., L. Hail, and C. Leuz. 2008. ``Mandatory IFRS Reporting around the World: Early Evidence on the Economic Consequences.''  \textit{Journal of Accounting Research} 46: 1085-1142.

\noindent Holland, P. 1986. ``Statistics and Causal Inference.''  \textit{Journal of the American Statistical Association} 81: 945-960.

\noindent Hope, O., and W. Thomas. 2008. ``Managerial Empire Building and Firm Disclosure.'' \textit{Journal of Accounting Research} 46: 591-626.

\noindent Hosman, C., B. Hansen, and P. Holland. 2010. ``The Sensitivity of Linear Regression Coefficients' Confidence Limits to the Omission of a Confounder.''  \textit{Annals of Applied Statistics} 4: 849-870.

\noindent Huang, J., Y. Shen, and Q. Sun. 2011. ``Nonnegotiable Shares, Controlling Shareholders, and Dividend Payments in China.''  \textit{Journal of Corporate Finance} 17: 122-133.

\noindent Iliev, P. 2010. ``The Effect of SOX Section 404:  Costs, Earnings Quality, and Stock Prices.''  \textit{Journal of Finance} 65: 1163-1196.

\noindent Imai, K., and D. Van Dyk. 2004. ``Causal Inference with General Treatment Regimes: Generalizing the Propensity Score.''\textit{ Journal of the American Statistical Association }99: 854-866.

\noindent Imbens, G., and T. Lemieux. 2008. ``Regression Discontinuity Designs:  A Guide to Practice.'' \textit{Journal of Econometrics} 142: 615-635.

\noindent Imbens, G., and P. Rosenbaum. 2005. ``Robust, Accurate Confidence Intervals with a Weak Instrument:  Quarter of Birth and Education.'' \textit{Journal of the Royal Statistical Society A} 168: 109-126.

\noindent Imbens, G., and D. Rubin. 2015. ``Causal Inference in Statistics, Social, and Biomedical Sciences.''

\noindent Imbens, G. W., and J. M. Wooldridge.  2009.  ``Recent Developments in the Econometrics of Program Evaluation.''  \textit{Journal of Economic Literature} 47:  5-86

\noindent John, K., and L. Litov. 2010. ``Managerial Entrenchment and Capital Structure:  New Evidence.'' \textit{Journal of Empirical Legal Studies} 7: 693-742.

\noindent Johnson, S., P. Boone, A. Breach, and E. Friedman. 2000. ``Corporate Governance in the Asian Financial Crisis.''  \textit{Journal of Financial Economics} 58: 141-186.

\noindent Kamar, E., P. Karaca-Mandic, and E. Talley. 2008. ``Going-Private Decisions and the Sarbanes-Oxley Act of 2002: A Cross-Country Analysis.'' \textit{ Journal of Law, Economics, and Organization} 25: 107-133.

\noindent Karpoff, J. M., and M. D. Wittry.  2015.  ``Test identification with legal changes:  The case of state antitakeover laws.''  Working paper, at http://ssrn.com/abstract=2493913. 

\noindent Keele, L., and J. Morgan. 2013. ``Stronger Instruments by Design.'' working paper, at http://ssrn.com/abstract=2280347. \textbf{}

\noindent Kezdi, G. 2004. ``Robust Standard Error Estimation in Fixed-Effects Panel Models." \textit{Hungarian Statistical Review} 9: 95-116.

\noindent Kinney Jr., W., and M. Shepardson. 2011. ``Do Control Effectiveness Disclosures Require SOX 404(b) Internal Control Audits? A Natural Experiment with Small U.S. Public Companies.'' \textit{Journal of Accounting Research} 49: 413-448.

\noindent La Porta, R., F. Lopez-de-Silanes, and A. Shleifer. 2008. ``The Economic Consequences of Legal Origins.'' \textit{Journal of Economic Literature} 46: 285-332.

\noindent Larcker, D., G. Ormazabal, and D. Taylor. 2011. ``The Market Reaction to Corporate Governance Regulation.''  \textit{Journal of Financial Economics }101: 431-448.

\noindent Larcker, D., and T. Rusticus. 2010. ``On the Use of Instrumental Variables in Accounting Research.'' \textit{Journal of Accounting and Economics} 49: 186-205.

\noindent Leamer, E. 1978. ``Specification Searches:  Ad Hoc Inference with Nonexperimental Data.''

\noindent Lee, D., and T. Lemieux. 2010. ``Regression Discontinuity Designs in Economics.''  \textit{Journal of Economic Literature} 48: 281-355.

\noindent Lee, D. and K. Suh Park. 2009. ``Does Institutional Activism Increase Shareholder Wealth? Evidence from Spillovers on Non-Target Companies.'' \textit{Journal of Corporate Finance} 15: 488-504.

\noindent Lemieux, T., and K. Milligan. 2008. ``Incentive Effects of Social Assistance:  A Regression Discontinuity Approach.''  \textit{Journal of Econometrics} 142: 807-828.

\noindent Lemmon, M. and K. Lins. 2003. ``Ownership Structure, Corporate Governance, and Firm Value: Evidence from the East Asian Financial Crisis.''  \textit{Journal of Finance} 58: 1445-1468.

\noindent Lennox, C., J. Francis, and Z. Wang. 2012. ``Selection Models in Accounting Research.'' \textit{The Accounting Review} 87: 589-616.

\noindent Li, H., M. Pincus, and S. Olhoft Rego. 2008. ``Market Reaction to Events Surrounding the Sarbanes-Oxley Act of 2002 and Earnings Management.''  \textit{Journal of Law and Economics} 51: 111-134.

\noindent Lin, C., and D. Su. 2008.  ``Industrial Diversification, Partial Privatization, and Firm Valuation:  Evidence from Publicly Listed Firms in China.''  \textit{Journal of Corporate Finance} 14:  405-417.

\noindent Litvak, K. 2007. ``The Effect of the Sarbanes-Oxley Act on Non-US Companies Cross-Listed in the US.'' \textit{Journal of Corporate Finance} 13: 195-228.

\noindent Litvak, K. 2008. ``Long-Term Effect of Sarbanes-Oxley on Cross-Listing Premia.'' \textit{European Financial Management} 14: 875-920.

\noindent Litvak, K. 2014. ``The Effect of Jurisdictional Competition for Corporate Charters on Financial and Accounting Measures of Firm Performance.'' working paper, at http://ssrn.com/abstract=1990805. 

\noindent Lo, K. 2003. ``Economic Consequences of Regulated Changes in Disclosure: the Case of Executive Compensation.''  \textit{Journal of Accounting and Economics} 35: 285-314.

\noindent Low, A. 2009. ``Managerial Risk-Taking Behavior and Equity-Based Compensation.'' \textit{Journal of Financial Economics} 92: 470-490.

\noindent MacKinlay, A. 1997. ``Event Studies in Economics and Finance.'' \textit{Journal of Economic Literature} 35: 13-39.

\noindent MacKinnon J. G, and M. D. Webb.  2014.  ``Wild Bootstrap Inference for Wildly Different Cluster Sizes.'' working paper, at http://hdl.handle.net/10419/97471.

\noindent Malani, A., and J. Reif.  2015.  ``Interpreting Pre-trends as Anticipation:  Impact on Estimated Treatment Effects from Tort Reform.''  \textit{Journal of Public Economics} forthcoming.

\noindent Malmendier, U., and G. Tate. 2009. ``Superstar CEOs.'' \textit{Quarterly Journal of Economics} 124: 1593-1638.

\noindent McCrary, J. 2008. ``Manipulation of the Running Variable in the Regression Discontinuity Design:  A Density Test.'' \textit{Journal of Econometrics} 142: 698-714.

\noindent Meng, R., X. Ning, X. Zhou, and H. Zhu. 2011. ``Do ESOPs Enhance Firm Performance? Evidence from China's Reform Experiment.''  \textit{Journal of Banking and Finance} 35: 1541-1551.

\noindent Mitchell, M., and J. Netter. 1989. ``Triggering the 1987 Stock Market Crash:  Antitakeover Provisions in the Proposed House Ways and Means Tax Bill?''  \textit{Journal of Financial Economics} 24: 37-68.

\noindent Morck, R., D. Yavuz, and B. Yeung. 2011. ``Banking System Control, Capital Allocation, and Economic Performance.''  \textit{Journal of Financial Economics} 100: 264-283.

\noindent Morgan, S., and C. Winship. 2014. ``Counterfactuals and Causal Inference:  Methods and Principles for Social Research.'' second edition (Cambridge Univ. Press).

\noindent Mulherin, H. 2007. ``Measuring the costs and benefits of regulation: Conceptual issues in securities markets.'' \textit{Journal of Corporate Finance} 13: 421-437.

\noindent Muravyev, A. 2013. ``Investor Protection and the Value of Shares:  Evidence from Statutory Rules Governing Variations of Shareholders' Class Rights in an Emerging Market.'' \textit{Journal of Law, Economics, and Organization} 29: 1344-1383. 

\noindent Murphy, K., and T. Sandino. 2010. ``Executive Pay and ``Independent'' Compensation Consultants.'' \textit{Journal of Accounting and Economics} 49: 247-262.

\noindent Nagar, V., D. Nanda, and P. Wysocki. 2003. ``Discretionary Disclosure and Stock-Based Incentives.''  \textit{Journal of Accounting and Economics} 34: 283-309.

\noindent Nasev, J., B. Black, and W. Kim.  2014.  ``Does Corporate Governance Affect Firms' Financial Reporting and Performance?  Panel Data versus Shock-Based Methods.''  Working paper, at http://ssrn.com/abstract=2133283.

\noindent Nenova, T. 2005. ``Control Values and Changes in Corporate Law in Brazil.''  \textit{Latin American Business Review} 6: 1-37.

\noindent Nguyen, B.D. and K.M. Nielsen. 2010. ``The Value of Independent Directors: Evidence from Sudden Deaths.''  \textit{Journal of Financial Economics} 98: 550-567. 

\noindent Oster, E. 2013. ``Unobservable Selection and Coefficient Stability:  Theory and Validation in Public Health.'' working paper, at http://ssrn.com/abstract=2266720.\textbf{}

\noindent Paik, M., B. Black, and D. Hyman. 2013. ``The Receding Tide of Medical Malpractice Litigation Part 2:  Effect of Damage Caps.'' \textit{Journal of Empirical Legal Studies} 10: 639-670.

\noindent Paik, M., B. Black, and D.A. Hyman.  2015a. ``Damage Caps and the Labor Supply of Physicians:  Evidence from the Third Reform Wave.''  Working paper, at http://ssrn.com/abstract=2470370.

\noindent Paik, M., B. Black, and D.A. Hyman.  2015b. ``Damage Caps and Defensive Medicine, Revisited.''  Working paper, at http://ssrn.com/abstract=2110656. 

\noindent Perez-Gonzalez, F. 2006. ``Inherited Control and Firm Performance.''  \textit{American Economic Review} 96: 1559-1588.

\noindent Petersen, M. 2009. ``Estimating Standard Errors in Finance Panel Data Sets:  Comparing Approaches.''  \textit{Review of Financial Studies} 22: 435-480.

\noindent Qiu, J., and F. Yu. 2009. ``The Market for Corporate Control and the Cost of Debt.''  \textit{Journal of Financial Economics} 93: 505-524.

\noindent Rauh, J. 2006. ``Own Company Stock in Defined Contribution Pension Plans:  A Takeover Defense?''  \textit{Journal of Financial Economics} 81: 379-410.

\noindent Roberts, M., and T. Whited. 2013. ``Endogeneity in Empirical Corporate Finance, in George M. Constantinides, Milton Harris, and Rene M. Stulz.'' eds., \textit{Handbook of the Economics of Finance}, vol. 2A, 493-572.

\noindent Rosenbaum, P. 2009. ``Design of Observational Studies.''

\noindent Rosenzweig, M., and K. Wolpin. 2000. ``Natural `Natural Experiments' in Economics.'' \textit{Journal of Economic Literature} 38: 827-874.

\noindent Rubin, D. 1974. ``Estimating Causal Effects of Treatments in Randomized and Nonrandomized Studies.''  \textit{Journal of Educational Psychology} 66: 688-701 (1974).

\noindent Rubin, D. 2008. ``For Objective Causal Inference, Design Trumps Analysis.'' \textit{Annals of Applied Statistics}, 2: 808-840.

\noindent Ryan Jr., H., and R. Wiggins III. 2004. ``Who is in whose pocket? Director compensation, board independence, and barriers to effective monitoring.''  \textit{Journal of Financial Economics} 73: 497-524.

\noindent Salas, J. 2010. ``Entrenchment, Governance, and the Stock Price Reaction to Sudden Executive Deaths.''  \textit{Journal Banking and Finance} 34: 656-666.

\noindent Shadish, W., T. Cook, and D. Campbell. 2002. ``Experimental and Quasi-Experimental Designs for Generalized Causal Inference.''

\noindent Shanken, J., and A. Tamayo. 2012. ``Payout Yield, Risk, and Mispricing:  A Bayesian Analysis.''  \textit{Journal of Financial Economics} 105: 131-152.

\noindent Small, D., and P. Rosenbaum. 2008. ``War and Wages:  The Strength of Instrumental Variables and their Sensitivity to Unobserved Biases.''  \textit{Journal of the American Statistical Association} 103: 924-933.

\noindent Smith, J., and P. Todd. 2005. ``Does Matching Overcome LaLonde's Critique of Nonexperimental Estimators.''  \textit{Journal of Econometrics} 125: 305-353.

\noindent Stock, J., J. Wright, and M. Yogo. 2002. ``A Survey of Weak Instruments and Weak Identification in Generalized Method of Moments.''  \textit{Journal of Business and Economic Statistics} 20: 518-529.

\noindent Strebulaev, I., and T. Whited. 2013. ``Dynamic Models and Structural Estimation in Corporate Finance.''  \textit{Foundations and Trends in Finance} 6: 1-163.

\noindent Stuart, T., and S. Yim. 2010. ``Board Interlocks and the Propensity to Be Targeted in Private Equity Transactions.''  \textit{Journal of Financial Economics} 97: 174-189.

\noindent Subramanian, G. 2004. ``The Disappearing Delaware Effect.''  \textit{Journal of Law, Economics, and Organization} 20: 32-59.

\noindent Thompson, S. 2010. ``Simple Formulas for Standard Errors that Cluster by Both Firm and Time.'' \textit{Journal of Financial Economics} 99: 1-10.

\noindent Wang, X. 2010. ``Increased Disclosure Requirements and Corporate Governance Decisions: Evidence from Chief Financial Officers in the Pre- and Post--Sarbanes-Oxley Periods.''  \textit{Journal of} \textit{Accounting Research} 48: 885-920.

\noindent Welch, I. 2012. ``A Critique of Recent Quantitative and Deep-Structure Modeling in Capital Structure Research and Beyond.''  \textit{Critical Finance Review} 2: 131-172. 

\noindent Wintoki, M.B. 2007. ``Corporate boards and regulation: The Effect of the Sarbanes--Oxley Act and the Exchange Listing Requirements on Firm Value.''  \textit{Journal of Corporate Finance} 13: 229-250. 

\noindent Wooldridge, J. 2010. \textit{Econometric Analysis of Cross Section and Panel Data}, 2nd edition.

\noindent Wooldridge, J. 2012. \textit{Introductory Econometrics:  A Modern Approach, }5th edition.

\noindent Yun, H. 2008. ``The Choice of Corporate Liquidity and Corporate Governance.'' \textit{ Review of Financial Studies} 22: 1447-1475.

\noindent Zhang, I.X. 2007. ``Economic Consequences of the Sarbanes--Oxley Act of 2002.'' \textit{Journal of Accounting and Economics} 44: 74-115.
